\section{Part (a): Steady-state temperature field}\label{sec:a}

The source line lies between grid nodes, so the main question is how it enters the discrete equations.
Section~\ref{sec:lines} compares six treatments on the given grid ($h=1$). A refined-grid reference $T^*$ and a grid-free
Monte Carlo estimate then measure their error. All numbers come from \texttt{code/part\_a/run\_part\_a.m}.

\textbf{Answer in brief.} Table~\ref{tab:answer} gives the nine nodal temperatures on the given grid (bold, treatment M3) and
the mesh-converged reference $T^*$ (brackets), which is the better estimate of the true field. M3 is within $3.1$~K of $T^*$,
which is $10.2\%$ of the local temperature rise above $200$~K. That closeness is specific to this geometry (Table~\ref{tab:robust}),
so the difference to $T^*$ is the uncertainty to carry into parts (b) and (c). The reasoning is in Sections~\ref{sec:lines} to
\ref{sec:validation} and the conclusion in Section~\ref{sec:result}.

\begin{table}[H]
\centering\small
\input{\rp tables/a_answer}
\caption{Steady-state temperatures [K] at the nine interior nodes. Bold: M3 on the given grid ($h=1$).
Brackets: reference $T^*$ ($\pm0.5$~K). Largest deviation $3.1$~K.}
\label{tab:answer}
\end{table}

% ======================================================================
\subsection{Finite-difference equations}

Taylor expansion about a node with neighbours at distance $a$ (left) and $c$ (right) gives
\begin{gather}
  T''\approx\frac{2}{a+c}\left[\frac{T_+-T}{c}-\frac{T-T_-}{a}\right]
  = w_-T_- - (w_-+w_+)\,T + w_+T_+, \label{eq:fd}\\
  \text{with}\quad w_-=\frac{2}{a(a+c)},\quad w_+=\frac{2}{c(a+c)}. \nonumber
\end{gather}
The right-hand side equals $T''+\tfrac{c-a}{3}T'''+\dots$, so for $a=c=h$ it is the central difference with error $O(h^2)$,
and for $a\ne c$ the error is first order. For any $a,c$ the formula is exact for quadratics.

Applying \eqref{eq:fd} in $x$ and $y$ turns $T_{xx}+T_{yy}=0$ into a linear system $A\mathbf T=\mathbf b$.
Away from the line the stencil is the usual 5-point one:
$\T{i-1}{j}+\T{i+1}{j}+\T{i}{j-1}+\T{i}{j+1}-4\T ij=0$. The conductivity $k$ does not appear: for uniform $k$ the temperature
field is independent of it, so $k$ matters only for parts (b) and (c).

% ======================================================================
\subsection{Six ways to use the line}\label{sec:lines}

Nodes $(1,j)$ have the line $0.5$ to their east and nodes $(2,j)$ have it $0.5$ to their west ($j=2,3$).

\textbf{M1, Shortley-Weller stencil.} Use \eqref{eq:fd} with $a=1$, $c=0.5$ and the known line value $\Ts$:
\begin{equation}
  T_{xx}\big|_{(1,j)}\approx \tfrac43\,T_{0,j}-4\,\T 1j+\tfrac83\,\Ts ,\qquad
  T_{xx}\big|_{(2,j)}\approx \tfrac83\,\Ts-4\,\T 2j+\tfrac43\,\T 3j .
\end{equation}
M1 is mathematically a quadratic ghost node: the parabola through the node, the line and the opposite neighbour gives
$T^g_{2,j}=\tfrac13T_{0,j}-2\T 1j+\tfrac83\Ts$, and $T_{0,j}-2\T 1j+T^g_{2,j}$ has the weights above.

\textbf{M2, linear ghost node.} Insert a temporary node where the missing neighbour would be, for example $(2,j)$ seen from
$(1,j)$, and fill it from the straight line through the node and the source line, $T^g=\sum_k L_k(x_g)\,T_k$.
The ghost lies beyond the line, so this is an extrapolation:
\begin{equation}
  T^g_{2,j}= L_1(2)\,\T 1j + L_{1.5}(2)\,\Ts = -\T 1j + 2\Ts .
\end{equation}
The 5-point stencil is then used with $T^g$ in place of the missing neighbour.

\textbf{M3, cubic ghost node.} Same ghost node, filled by the cubic through $x=0,\,1,\,1.5,\,3$
(opposite neighbour, node, line, and the node beyond the ghost), which is an interpolation:
\begin{equation}
  T^g_{2,j}=\tfrac19\,T_{0,j}-\T 1j+\tfrac{16}{9}\,\Ts+\tfrac19\,\T 3j .
\end{equation}
This gives $T_{xx}|_{(1,j)}\approx\tfrac{10}{9}T_{0,j}-3\,\T 1j+\tfrac19\T 3j+\tfrac{16}{9}\Ts$, with a mirror image for $(2,j)$.
The linear fit is exact for straight lines, the parabola of M1 for quadratics, and this cubic for cubics, so its local error
in $x$ is $O(h^2)$ instead of $O(h)$.

\textbf{M4 to M6, snapping.} Naive baselines that ignore the half-cell offset: clamp the nearest nodes to $350$~K.
M4 clamps $(1,2),(1,3)$, M5 clamps $(2,2),(2,3)$, M6 clamps all four. They show what is lost by not modelling the offset.

\medskip
\textbf{The nine equations for M3} ($\Tb=200$, $\Ts=350$, ordered by node), carried through the solver, time-marching and
random-walk sections:
\begin{align}
(1,1)\!:\ & \T 21+\T 12-4\T 11 = -400, & (2,1)\!:\ & \T 11+\T 31+\T 22-4\T 21 = -200, \nonumber\\
(3,1)\!:\ & \T 21+\T 32-4\T 31 = -400, & (1,2)\!:\ & \T 11+\T 13+\tfrac19\T 32-5\T 12 = -\tfrac{7600}{9}, \nonumber\\
(2,2)\!:\ & \T 21+\T 23+\tfrac{10}{9}\T 32-5\T 22 = -\tfrac{5800}{9}, & (3,2)\!:\ & \T 31+\T 33+\T 22-4\T 32 = -200, \nonumber\\
(1,3)\!:\ & \T 12+\tfrac19\T 33-5\T 13 = -\tfrac{9400}{9}, & (2,3)\!:\ & \T 22+\tfrac{10}{9}\T 33-5\T 23 = -\tfrac{7600}{9}, \nonumber\\
(3,3)\!:\ & \T 23+\T 32-4\T 33 = -400. & & \label{eq:system}
\end{align}
Every row is strictly diagonally dominant with non-negative off-diagonal entries, so Jacobi and Gauss-Seidel converge.

\begin{table}[H]
\centering\small
\input{\rp tables/a_treatments}
\caption{Temperatures [K] at the nine nodes (Gauss elimination). M1 Shortley-Weller, M2 linear ghost, M3 cubic ghost,
M4 to M6 snapping (left pair, right pair, all four). $T^*$ is the mesh-converged reference of Section~\ref{sec:refine}.
``max dev.'' is the largest $|T-T^*|$ over the nine nodes, in kelvin and as a percentage of the local rise $T^*-200$~K.}
\label{tab:treatments}
\end{table}

\begin{figure}[H]
  \centering
  \includegraphics[width=0.88\textwidth]{a_treatments.pdf}
  \caption{Nodal temperature maps for the six treatments (dark red = clamped to $350$~K).}
  \label{fig:treatments}
\end{figure}

\textbf{Is the ranking general?} The $h=1$ systems only see that the line passes the node rows $y=2$ and $3$, so they are
identical for any line that does. The converged field is not. Table~\ref{tab:robust} rebuilds $T^*$ for four such lines
(same procedure as Section~\ref{sec:refine}) and compares M1 to M3 with it.
\begin{table}[H]
\centering\small
\input{\rp tables/a_robust}
\caption{Maximum deviation [K] of M1 to M3 from the reference $T^*$ rebuilt for each line $x=1.5$, $y\in[y_{lo},y_{hi}]$.
The first row is the given problem. Bold: smallest in the row.}
\label{tab:robust}
\end{table}
M3 is closest only for the given line. For the two longer lines it is the worst of the three, and M1 stays between $13.5$ and
$17.2$~K throughout. The error at $h=1$ is therefore dominated by the line ends, which a grid with spacing $1$ cannot
resolve, and the better local accuracy of M3 helps only when those ends happen to be benign. M3 is reported for the given
problem because it is closest to $T^*$ here, as an empirical result rather than a general ranking.

% ======================================================================
\subsection{Solving \texorpdfstring{$A\mathbf T=\mathbf b$}{A T = b}}

All solvers are applied to the M3 system \eqref{eq:system}. The iterative solvers start from $\mathbf T=\mathbf 0$ and stop
when the largest residual falls below $10^{-8}$; one sweep is one pass over all nodes.

\textbf{Direct.} Gauss elimination with partial pivoting, and Crout LU ($A=LU$, then $L\mathbf y=\mathbf b$,
$U\mathbf x=\mathbf y$). They agree to $10^{-13}$~K.

\textbf{Iterative.} With relaxation factor $\omega$,
\begin{equation}
  \text{Jacobi: } x_i^{(k+1)}=\frac{b_i-\sum_{j\ne i}A_{ij}x_j^{(k)}}{A_{ii}},\qquad
  \text{SOR: } x_i^{\rm new}=(1-\omega)x_i^{\rm old}+\omega\,\frac{b_i-\sum_{j\ne i}A_{ij}x_j}{A_{ii}}.
\end{equation}
SOR uses the newest values immediately, and $\omega=1$ is Gauss-Seidel (Liebmann's method).

\begin{table}[H]
\centering\small
\input{\rp tables/a_solvers}
\caption{Solvers on the $9\times9$ system of M3. The residual factor is the mean reduction of $\|\mathbf b-A\mathbf x\|_\infty$
per sweep over the second half of the iteration history.}
\label{tab:solvers}
\end{table}

\begin{figure}[H]
  \centering
  \includegraphics[width=0.85\textwidth]{a_solvers.pdf}
  \caption{Left: residual history. Right: sweeps against $\omega$ (minimum at $\omega=1.15$, $16$ sweeps at $1.10$).}
  \label{fig:solvers}
\end{figure}

\begin{itemize}
  \item Gauss-Seidel converges at about the square of the Jacobi rate ($0.346$ against $0.586^2=0.343$). Theory gives exactly
        the square for the plain 5-point matrix; the extra $\tfrac19$ couplings of M3 change it slightly.
  \item SOR with $\omega=1.15$ needs $15$ sweeps against $25$ for Gauss-Seidel. The scan is flat between $1.10$ and $1.15$
        (theory from the Jacobi rate gives $1.10$), and the cost rises steeply towards $\omega=2$ ($538$ sweeps at $1.95$).
  \item The gap grows with grid size. On the grids of Section~\ref{sec:refine} ($N=8,16,32$) Gauss-Seidel needs
        $103,\,407,\,1562$ sweeps ($\propto N^2$) and SOR with $\omega=2/(1+\sin(\pi/N))$ needs $38,\,74,\,151$ ($\propto N$).
\end{itemize}

% ======================================================================
\subsection{Marching in time to the steady state}

The steady solution is the limit of $\mathrm dT/\mathrm ds=\nabla^2T$ as pseudo-time $s\to\infty$ (diffusivity $1$; $s$ is not
the physical time of part (c)). After spatial discretisation,
\begin{equation}
  \frac{\mathrm d\mathbf T}{\mathrm ds}=A\mathbf T-\mathbf b,\qquad \mathbf T(0)=200\ \mathrm{K},\qquad
  \mathbf T(s)=\mathbf T_{ss}+e^{As}\left(\mathbf T(0)-\mathbf T_{ss}\right).
\end{equation}
The last form is the exact solution, used as the reference ($e^{As}$ by a Taylor series with scaling and squaring).
Power and inverse iteration give $\lambda_{\max}=-6.98$ and $\lambda_{\min}=-1.73$, so the slowest time constant is
$\tau=1/|\lambda_{\min}|=0.577$. Explicit Euler is stable only for
\begin{equation}
  \Delta s\le\frac{2}{|\lambda_{\max}|}=0.2864,\qquad\text{and RK4 for }\Delta s\le \frac{2.785}{|\lambda_{\max}|}=0.399.
\end{equation}
Running explicit Euler at $0.95$ and $1.05$ of its limit confirms the eigenvalue: the error falls to $4\times10^{-7}$~K in the
first case and grows to $7\times10^{5}$~K in the second (Fig.~\ref{fig:transient}, right).

\begin{table}[H]
\centering\small
\input{\rp tables/a_transient}
\caption{Time-marching schemes to $s=10\tau$ against the exact solution. All end at the steady state of M3. The last column is
the error as a percentage of the rise of $T_{2,2}$ above $200$~K. The last row has the same work as RK4 ($4\times$ the steps).}
\label{tab:transient}
\end{table}

\begin{figure}[H]
  \centering
  \includegraphics[width=0.9\textwidth]{a_transient.pdf}
  \caption{Left: $T_{2,2}(s)$ for several schemes, a marker at every step. Right: explicit Euler just below and just above
  its stability limit.}
  \label{fig:transient}
\end{figure}

RK4 is about $170$ times more accurate than explicit Euler at the same step ($0.05$ against $8.9$~K) and about $30$ times at
equal work (Euler at a quarter of the step, $1.7$~K). Implicit Euler is unconditionally stable but strongly damped.
Crank-Nicolson is second order ($0.64$~K at $0.5\,\Delta s_{\max}$) but rings at $10\,\Delta s_{\max}$ (error $48$~K): its
amplification factor $(1+z/2)/(1-z/2)$ tends to $-1$ for fast modes, so stability does not guarantee a smooth transient.
Settling to $0.1$~K takes $s=3.61\approx6.3\tau$. For a steady answer alone, time marching costs far more than one elimination.

% ======================================================================
\subsection{Grid refinement and Richardson extrapolation}\label{sec:refine}

On a uniform grid with $N=8,16,\dots,128$ divisions the line falls exactly on nodes, so it is imposed as fixed nodes and no
special stencil is needed. These $N^2$-unknown systems are solved by matrix-free SOR with $\omega=2/(1+\sin(\pi/N))$.
The observed order $p$ comes from the three finest grids, and Richardson extrapolation gives the reference $T^*$:
\begin{equation}
  p=\log_2\frac{T_h-T_{h/2}}{T_{h/2}-T_{h/4}},\qquad
  T^*=T_{h/4}+\frac{T_{h/4}-T_{h/2}}{2^{p}-1}.
\end{equation}

\begin{table}[H]
\centering\small
\input{\rp tables/a_refinement}
\caption{Refinement sequence at three nodes [K], observed order $p$, and the extrapolated reference $T^*$.}
\label{tab:refine}
\end{table}

\begin{figure}[H]
  \centering
  \begin{subfigure}[t]{0.47\textwidth}\centering
    \includegraphics[width=\textwidth]{a_convergence.pdf}
    \caption{Change between successive grids, all nine nodes.}\label{fig:conv}
  \end{subfigure}\hfill
  \begin{subfigure}[t]{0.47\textwidth}\centering
    \includegraphics[width=\textwidth]{a_field.pdf}
    \caption{Converged field ($N=128$). Dots are the given nodes.}\label{fig:field}
  \end{subfigure}
  \caption{Grid convergence and the converged temperature field.}
\end{figure}

The observed order is $p\approx1.0$ at every node, not the $2$ expected for a smooth problem. Fig.~\ref{fig:conv} shows it
without using $T^*$: halving $h$ halves the change between grids. At the ends of a fixed-temperature segment the solution
behaves like $\sqrt r$, so $T_{xx}\propto r^{-3/2}$ is unbounded and the Taylor argument behind $O(h^2)$ fails.
The exact slit-tip solution of Section~\ref{sec:validation} shows the same order $1$ at fixed points, so this comes from the
geometry and not from a bug.

Assuming $p=2$ would shift $T^*$ by about $0.4$~K, so the measured $p$ is used. Repeating the extrapolation one level
finer ($N=64,128,256$) changes $T^*$ by at most $0.03$~K, and the grid-free walk-on-spheres estimate below agrees within
$0.3$~K. The bound $\pm0.5$~K is a judgement from these three checks and not a proven error estimate.

% ======================================================================
\subsection{Monte Carlo random walks (beyond the syllabus)}

Laplace's equation has the mean-value property, so a temperature equals the expected boundary temperature reached by a
random walker started there (Kakutani, 1944). This is the one method used that is not in the lecture notes. Two variants:

\textbf{Grid walk.} Row $r$ of $A\mathbf T=\mathbf b$ is rewritten as $T_r=\sum_kp_{rk}T_k+q_r\bar T_r$ with
$p_{rk}=A_{rk}/(-A_{rr})$ (step probabilities) and $q_r=1-\sum_kp_{rk}$ (probability of hitting the edge or line).
A walker steps with probabilities $p_{rk}$ and scores $\bar T_r$ when it hits. The mean score reproduces the finite-difference
answer, so it checks the solve and not the model. This needs $p_{rk}\ge0$, which holds for M3.

\textbf{Walk on spheres.} No grid. From the current point, jump to a random point on the largest circle that touches the plate
edge or the source segment. Repeat until within $\varepsilon=10^{-4}$ of either, then score $350$ or $200$~K. The distance to
the segment is exact, so the half-cell offset never appears. The remaining errors are the statistical error ($\propto N^{-1/2}$)
and the thin absorption layer $\varepsilon$.

\begin{table}[H]
\centering\small
\input{\rp tables/a_montecarlo}
\caption{Monte Carlo estimates [K] with $2\times10^5$ walks per node (fixed seed), with 95\% half-widths.}
\label{tab:mc}
\end{table}

\begin{figure}[H]
  \centering
  \includegraphics[width=0.46\textwidth]{a_montecarlo.pdf}
  \caption{Rms error over the nine nodes (5 repeats) against the number of walks $N$. The slope is $-1/2$.}
  \label{fig:mc}
\end{figure}

The grid walk matches M3 to within $0.28$~K at every node, inside its 95\% interval at all nine. Walk-on-spheres lands within
$0.3$~K of $T^*$ at all nine nodes, so it confirms the Richardson reference independently of the grid. A tenfold gain in
accuracy needs a hundredfold more walks, so Monte Carlo is a check rather than a competitive solver. Both variants are also
tested against exact answers in Section~\ref{sec:validation}.

% ======================================================================
\subsection{Validation of the scripts}\label{sec:validation}

\begin{table}[H]
\centering\small
\begin{tabular}{@{}P{2.6cm} P{2.6cm} P{5.4cm} P{4.2cm}@{}}
\toprule
Test & Script checked & What is checked & Result \\
\midrule
Lecture exercise & \texttt{sor} & $2\times2$ interior, $\omega=1.1$: $46.9,\ 71.9,\ 40.6,\ 65.6$ & matches to $0.05$~K \\
Smooth solution & \texttt{sorGrid} & $u=e^x\sin y$ on the unit square (Table~\ref{tab:test}) & order $2.00$ \\
Quadratic patch & \texttt{buildSystem} & $u=x^2-y^2+xy$ on edges and line: M1, M3 exact & $1.8\times10^{-15}$~K; M2 $0.15$~K \\
Cubic patch & \texttt{buildSystem} & $u=x^3-3xy^2$: only M3 should be exact & M3 $3.6\times10^{-15}$~K; M1 $0.23$, M2 $0.68$~K \\
Slit tip & \texttt{sorGrid} & $u=\mathrm{Re}\sqrt{z-z_0}$, zero on the slit (Table~\ref{tab:test}) & order $1.0$ at fixed points \\
Solver agreement & \texttt{gaussElim}, \texttt{luSolve}, \texttt{jacobi}, \texttt{sor} & against MATLAB's \texttt{\textbackslash}, used only here as an oracle & $<10^{-8}$ \\
Time marching & \texttt{marchHeat}, \texttt{expmSeries} & error ratio when $\Delta s$ is halved, against $e^{As}$ & Euler $2.0$, implicit Euler $2.0$, CN $4.0$, RK4 $17$ \\
Monte Carlo & \texttt{walkGrid}, \texttt{walkOnSpheres} & three nodes against the exact quadratic, within 4 standard errors & pass \\
Physics & \texttt{buildSystem} & $200\le T\le350$~K for M1 to M3 (maximum principle) & holds \\
\bottomrule
\end{tabular}
\caption{Validation tests (\texttt{tests/test\_part\_a.m}). All pass.}
\label{tab:validation}
\end{table}

\begin{table}[H]
\centering\small
\input{\rp tables/a_test}
\caption{Convergence on exact solutions. The smooth case is second order. The slit-tip solution has a $\sqrt r$ singularity like
the segment ends, and the error at five fixed points converges at first order, which is the behaviour seen for the plate in
Section~\ref{sec:refine}.}
\label{tab:test}
\end{table}

% ======================================================================
\subsection{Result and conclusion for part (a)}\label{sec:result}

\textbf{Question.} Steady-state temperature at the nine interior nodes of the given grid, from $T_{xx}+T_{yy}=0$.

\textbf{Equation used.} At every node, with the line entering through the ghost value $T^g$ (M3),
\begin{equation}
  \T{i-1}{j}+\T{i+1}{j}+\T i{j-1}+\T i{j+1}-4\T ij=0,\qquad
  T^g_{2,j}=\tfrac19\,T_{0,j}-\T 1j+\tfrac{16}{9}\,\Ts+\tfrac19\,\T 3j ,
\end{equation}
with a mirror image for the nodes $(2,j)$. This is assembled as $A\mathbf T=\mathbf b$, equation~\eqref{eq:system}, and solved
by Gauss elimination. The values are in Table~\ref{tab:answer}.

\textbf{Do the results change, and why?} Yes. The line is half a cell from four nodes, so how it enters the equations matters
more than the solver (all solvers agree to $10^{-8}$~K). Errors below are against $T^*$, as a percentage of the local rise
above $200$~K.
\begin{itemize}
  \item Snapping nodes to $350$~K moves the source by half a cell. Errors reach $84$~K, more than the rise itself.
  \item Shortley-Weller (M1) has a first-order local error $\tfrac{c-a}{3}T'''$ at the four nodes next to the line:
        $17.2$~K ($30.4\%$).
  \item The linear ghost node (M2) is closer ($11.6$~K, $22.5\%$), partly by error cancellation, and it fails the patch tests.
  \item The cubic ghost node (M3) is exact for cubics and closest on this problem: $3.1$~K ($10.2\%$). Table~\ref{tab:robust}
        shows this does not carry over to other line lengths.
\end{itemize}
Refinement converges at only first order because of the singular segment ends. Walk-on-spheres, which uses no grid, confirms
$T^*$ to within $0.3$~K.

\textbf{Most appropriate methods.}
\begin{itemize}
  \item Line treatment: M3 for the nodal values on the given grid, and $T^*$ from refinement as the benchmark.
        M1 is the steadier alternative if the line geometry changes.
  \item Linear solve: Gauss elimination for nine unknowns. For fine grids SOR with the optimal $\omega$, whose cost grows as $N$
        against $N^2$ for Gauss-Seidel.
  \item Time marching and Monte Carlo are cross-checks, not recommended for the answer: the first costs far more than one
        elimination, the second converges as $N^{-1/2}$.
\end{itemize}

\textbf{For parts (b) and (c).} Use the nodal values of Table~\ref{tab:answer} (file column \texttt{T\_cubic\_h1}; \texttt{T\_reference}
holds $T^*$). The conductivity does not enter $T$. The deviation from $T^*$ (up to $3.1$~K) is the temperature uncertainty to
carry forward, and a gradient divides temperature differences by a spacing of $0.5$ to $1$, so the uncertainty in $\vec q$ is
larger still, particularly near the line.

Equation~\eqref{eq:fd} is a second-derivative formula. A first derivative with neighbours at $a$ (left) and $c$ (right) needs
\begin{equation}
  T'\approx -\frac{c}{a(a+c)}\,T_- + \frac{c-a}{ac}\,T + \frac{a}{c(a+c)}\,T_+ ,
\end{equation}
which is exact for quadratics and reduces to the central difference for $a=c$. Next to the line, $a=1$ and $c=0.5$ (or the
reverse) in $x$.
